% Generated draft: requires mathematical and source review.
\documentclass[11pt]{article}
\usepackage[margin=25mm]{geometry}
\usepackage{fontspec}
\setmainfont{Latin Modern Roman}
\usepackage{amsmath,amssymb,mathtools,unicode-math}
\setmathfont{Latin Modern Math}
\usepackage{enumitem,needspace,xurl,hyperref}
\setlength{\parindent}{0pt}
\setlength{\parskip}{5pt}
\setlength{\emergencystretch}{3em}
\newcommand{\N}{\mathbb{N}}
\newcommand{\Z}{\mathbb{Z}}
\newcommand{\Q}{\mathbb{Q}}
\newcommand{\R}{\mathbb{R}}
\newcommand{\C}{\mathbb{C}}
\newcommand{\F}{\mathbb{F}}
\newcommand{\A}{\mathbb{A}}
\newcommand{\PP}{\mathbb{P}}
\newcommand{\E}{\mathbb{E}}
\newcommand{\eps}{\varepsilon}
\newcommand{\dd}{\,\mathrm{d}}
\begin{document}
\section*{UnsolvedMath ID: 11000121}
\textbf{Source problem code:} catalogue entry
\subsection{Definitions and mathematical statement}
For an integer $a\ge1$, let $N_a=\#^a\mathbb{RP}^2$ be the closed connected nonorientable surface of genus $a$, and let $\operatorname{Mod}(N_a)$ be its group of self-homeomorphisms modulo isotopy. There is no orientation-preserving restriction on a nonorientable surface.

\textbf{Problem 3.3.} For every pair of integers $g>h\ge1$, is the image of every group homomorphism
\[
 \varphi:\operatorname{Mod}(N_g)\longrightarrow\operatorname{Mod}(N_h)
\]
finite? In quantifiers, the proposed assertion is
\[
 \forall g>h\ge1\quad\forall\varphi\in\operatorname{Hom}(\operatorname{Mod}(N_g),\operatorname{Mod}(N_h)),
 \qquad |\varphi(\operatorname{Mod}(N_g))|<\infty.
\]
The source requires closed nonorientable surfaces and a genus drop, but places no injectivity, surjectivity, or continuity condition on $\varphi$ and no additional lower-genus threshold. The desired conclusion is finite image, not necessarily trivial image.

The research summary reports an orientable analogue and related rigidity results, but does not claim a resolution of this nonorientable statement. Curve-complex or automorphism rigidity alone is not substituted for the arbitrary-homomorphism assertion.
\subsection{Short English statement}
Must every homomorphism from the mapping class group of a closed nonorientable surface to one of lower nonorientable genus have finite image? The question includes all genus pairs stated in the source.
\subsection{Sources}
\begin{itemize}
\item Local UnsolvedMath registry, permanent ID 11000121, source code catalogue entry; no proposer or proposal year supplied.
\item Farb, edited volume \emph{Problems on Mapping Class Groups and Related Topics}, Problem 3.3, PDF page 98; \url{https://www.math.uchicago.edu/~farb/papers/mcgbook.pdf}. The registry records PDF-font-structure extraction and access on 2026-07-29; its inventory distinguishes 329 labeled headings by source page.
\item The supplied summary cites Harvey--Korkmaz, ``Homomorphisms from mapping class groups,'' \emph{Bulletin of the London Mathematical Society} 37 (2005), 275--284, for an orientable analogue with its own genus hypotheses. It also names Ivanov, Ivanov--McCarthy, and Atalan--Korkmaz for related rigidity work without full bibliographic entries.
\end{itemize}
\end{document}
